% Respaldo de la lamina 9: una estacion atacada (E1), PARAAD-SA con el umbral p99.9.
% Misma construccion que estacion_atacada.tex; solo cambian datos, color y escala.
% Datos: results/linea_base/rev24/metricas_banco.csv, filas PARAAD, nivel 0.001, E1.
% Banco 2024-jul 2026, 30 calendarios (10 por año); los conteos son la suma de los 30.
% VP = dia atacado con alarma, FP = dia normal con alarma, FN = dia atacado sin alarma.
% Unidad del eje Y: dias al año por estacion. Las coordenadas se escriben con los conteos
% del CSV (estacion-dias) y "y filter" las divide entre 630.27 estacion-años:
% 23 005 estacion-dias con dato / 365 = 63.03 años-estacion, por 10 calendarios por año.
% Por cada tamaño de ataque, dos barras apiladas que comparten el trozo VP:
%   izquierda = alarmas (VP + FP)   -> precision = VP / alto de la barra
%   derecha   = dias atacados (VP + FN) -> recall = VP / alto de la barra
% Truco de pgfplots: un solo eje "ybar stacked"; cada barra tiene su propia x
% (centro -0.24 y centro +0.24) y las tres series usan las mismas 12 x, con 0 donde
% la pieza no va (FP no va en la barra de atacados ni FN en la de alarmas).
% Posiciones: delta 20/50/100 % en 1, 2, 3; error fijo 10/20/40 ppb en 4.5, 5.5, 6.5.
% Los scope no funcionan dentro de este eje; por eso cada tipo de texto tiene un estilo.
\begin{tikzpicture}[
    etiqueta/.style={font=\sffamily\scriptsize, anchor=south, inner sep=2pt},
    tamano/.style={font=\sffamily\small, text=texto, anchor=north},
    familia/.style={font=\sffamily\scriptsize, text=tenue, anchor=north},
  ]
\begin{axis}[
    ejes tesis,
    ybar stacked,
    bar width=0.40,                 % sin unidad = unidades del eje x
    scale only axis, width=12.3cm, height=3.7cm,
    xmin=0.45, xmax=7.05,
    y filter/.expression={y/630.27},  % estacion-dias -> dias al año por estacion
    ymin=0, ymax=5.5, ytick={0,1,...,5},
    ylabel={Días al año por estación},
    ylabel style={font=\sffamily\small},
    xtick={0.76,1.24, 1.76,2.24, 2.76,3.24, 4.26,4.74, 5.26,5.74, 6.26,6.74},
    xticklabels={alarmas,atacados, alarmas,atacados, alarmas,atacados,
                 alarmas,atacados, alarmas,atacados, alarmas,atacados},
    xtick style={draw=none},
    xticklabel style={font=\sffamily\tiny, text=tenue},
    legend style={at={(0.5,1.02)}, anchor=south, legend columns=3,
                  /tikz/every even column/.append style={column sep=0.5cm}},
    legend image code/.code={\fill[#1] (0cm,-0.1cm) rectangle (0.3cm,0.1cm);},
    clip=false,
  ]
  % Orden de las series = orden en que se apilan (VP abajo).
  % x:                      d20-A  d20-D   d50-A  d50-D   d100-A d100-D  a10-A  a10-D   a20-A  a20-D   a40-A  a40-D
  \addplot[barra=paraad] coordinates {
    (0.76,58)   (1.24,58)   (1.76,343)  (2.24,343)  (2.76,910)  (3.24,910)
    (4.26,58)   (4.74,58)   (5.26,444)  (5.74,444)  (6.26,1402) (6.74,1402)};
  \addlegendentry{VP: ataque detectado}
  \addplot[barra=sinataque!70!white] coordinates {
    (0.76,1750) (1.24,0)    (1.76,1750) (2.24,0)    (2.76,1763) (3.24,0)
    (4.26,1751) (4.74,0)    (5.26,1751) (5.74,0)    (6.26,1758) (6.74,0)};
  \addlegendentry{FP: falsa alarma}
  \addplot[barra=paraad!30!white] coordinates {
    (0.76,0)    (1.24,1719) (1.76,0)    (2.24,1434) (2.76,0)    (3.24,867)
    (4.26,0)    (4.74,1719) (5.26,0)    (5.74,1333) (6.26,0)    (6.74,375)};
  \addlegendentry{FN: ataque no detectado}

  % Precision (sobre la barra de alarmas) y recall (sobre la de atacados); coinciden con el CSV.
  % "axis cs" no pasa por el y filter: la altura se divide aqui mismo entre 630.27.
  \node[etiqueta, text=tenue] at (axis cs:0.76,1808/630.27) {P\,0.03}; \node[etiqueta, text=paraad] at (axis cs:1.24,1777/630.27) {R\,0.03};
  \node[etiqueta, text=tenue] at (axis cs:1.76,2093/630.27) {P\,0.16}; \node[etiqueta, text=paraad] at (axis cs:2.24,1777/630.27) {R\,0.19};
  \node[etiqueta, text=tenue] at (axis cs:2.76,2673/630.27) {P\,0.34}; \node[etiqueta, text=paraad] at (axis cs:3.24,1777/630.27) {R\,0.51};
  \node[etiqueta, text=tenue] at (axis cs:4.26,1809/630.27) {P\,0.03}; \node[etiqueta, text=paraad] at (axis cs:4.74,1777/630.27) {R\,0.03};
  \node[etiqueta, text=tenue] at (axis cs:5.26,2195/630.27) {P\,0.20}; \node[etiqueta, text=paraad] at (axis cs:5.74,1777/630.27) {R\,0.25};
  \node[etiqueta, text=tenue] at (axis cs:6.26,3160/630.27) {P\,0.44}; \node[etiqueta, text=paraad] at (axis cs:6.74,1777/630.27) {R\,0.79};

  % Separador entre familias de ataque (como en la version anterior).
  \draw[tenue, densely dashed, line width=0.5pt] (axis cs:3.75,0) -- (axis cs:3.75,5.33);

  % Nombre de cada tamaño bajo sus dos barras y de cada familia debajo.
  \node[tamano] at ([yshift=-0.45cm]axis cs:1,0)   {$\delta$\,20\,\%};
  \node[tamano] at ([yshift=-0.45cm]axis cs:2,0)   {$\delta$\,50\,\%};
  \node[tamano] at ([yshift=-0.45cm]axis cs:3,0)   {$\delta$\,100\,\%};
  \node[tamano] at ([yshift=-0.45cm]axis cs:4.5,0) {10 ppb};
  \node[tamano] at ([yshift=-0.45cm]axis cs:5.5,0) {20 ppb};
  \node[tamano] at ([yshift=-0.45cm]axis cs:6.5,0) {40 ppb};
  \node[familia] at ([yshift=-0.95cm]axis cs:2,0)   {Anomalía proporcional};
  \node[familia] at ([yshift=-0.95cm]axis cs:5.5,0) {Escalón aditivo fijo};
\end{axis}
\end{tikzpicture}

